\documentclass{article}
\usepackage[utf8]{inputenc}
\usepackage{hyperref}
\usepackage{url}
\title{??? ??????? ???????? *Oryza sativa*(??)?? ????? 4?????(????--?????--???--??????) ?????? PFAS}
\author{sci-adk (deterministic render)}
\date{\today}

\begin{document}
\maketitle

\noindent\textit{Draft compiled by sci-adk from Spec \texttt{pfas-rice} (v1). Belief state is revisable as Evidence accrues.}

\section{Goal}
??? ??????? ???????? *Oryza sativa*(??)?? ????? 4?????(????--?????--???--??????) ?????? PFAS
??????????? ????????, DPU ????????????? ???????? ??????(IOC extension)???? ?????????? ???????.
???????????:

1. ?????????? ?????: PFAS??? ????? ???????????? ??????, ???? ???????? ???????? ????(GHK)??
   ??????? ????????(Michaelis--Menten)?? ????????????, ???? ???????? ?????? ??????$\cdot$????? ?????
   ????(advection)?? ????? ?????????? ??????????.
2. ????????$\cdot$??????????? ????????: ????? ?????(C4--C12)?? ????????(??????????? PFCA vs ?????????
   PFSA, ?????? ????????? GenX)?? ???? ?????????? ?????$\cdot$???????? ?????????, ??? ???????
   ????? ?????????? ??????????.
3. ?????--?????? ??????: ????? ???????? ????? C\_w\textasciicircum{}o(t)?? ???????? Q\_TP(t)??? ???? ?????
   ?????(HYDRUS-1D)???? ????? ?????? ODE??? ?????????? ?????????(Method A) ????????? ???????
   ??????????.
4. ?????? ????? ?????$\cdot$??????: ???????? ????????????? ?????? ??????????? ??????, ????????
   ?????/?????? ??????(Yamazaki, Tang, Kim)???? ??????? ?????(???????? ????, ?????? ??????,
   root>straw>grain?? ???????? ????????)??? ???????????.
5. ????????: ??????? PFAS ????(SMILES)??? ?????????? ????? ??????????????, ?????? ????????
   ???????????? ????????? ?? ?????? ??????? ??????????.

\section{Background}
???????????????(PFAS)??? ?????--????? ???????? ?????? ????????? ??????? "?????????
???????????(forever chemicals)"?? ???????, ?????????? ???? ????????? ????? ???????
?????????? ??????$\cdot$???$\cdot$?????? ????????? ??????????. ????? ???(paddy)??? PFAS ??????? ??????
???????. ????? ???? ????? ???????? ?????????, ?????--??????? ??????????? ????? ???????????
????????? ????? ???????? ????????(pore water) ???????? ???? ?????? ????? ???????? ??????????
????????? ????? ???????? ????????. ?????? "??? ??????? PFAS?? ?????(grain)?? ???????
????????????"??? ??????????????? ???????????? ??????????.

???????? ????? ??????????? ??????? PFAS?? ?????? ?????????? ??????????? ???????. Trapp/Briggs
??????? ???????? ??????????? ??????? DPU(Dynamic Plant Uptake) ?????????????? ??????
?????????????? ?????????--??? ?????????(Kow)?? ???? ?????? ????? ???????????? ????? ???? ??
?????. ?????? PFAS??? pKa?? ???? ????? ?????? pH???? ???????? 100\% ??????? ????????
(f\_d$\approx$1)???? ??????????. ??????????? ???????? ??????????? ??????????? ??????????????? ????????
(?????????? ???? e\textasciicircum{}N$\approx$107), ?????????????? Kow ????? ?????? ?????????? ????????. ??? PFAS??
???? ??????? ???? ???????? ????? ??????????? + ???????? ???? ????? + ?????????$\cdot$????????$\cdot$???????
?????? ????????????? ??????? ????????????? ???????????? ?????. ?????? ????? ????????? ??????????
???????? ???????? ??????? ?????? ????????????? ????????? ?????? ?????? ?????? ?????(terminal
accumulator)??????, ?????? BAF ??? ?????????? ???????? ?? ????? ?????????????? ?????? ???????
????????????? ?????????.

??????? ?????? ???????, ???????? ????????????(IOC)?? ????? ??????????, ??? ??????? ?????? ????????
?????????? ?? ????? ??????? ?????????.

\section{Method}
(1) ????? ?????????. ????????????? ?? ??????? ?????????? C\_k($\mu$g/kg)????, ??????(stiff) ODE???
BDF ??????? ???????????. ??? ???????? ???????????(?$\mu$$\cdot$C, $\mu$=M/M)?? (PFAS?? ????? ???? 0???)
?????????? ????????.

(2) ???? ????? j\_R(??????????). Goldman--Hodgkin--Katz(GHK) ???????????????? ?????? ????????
??????? ??????????? ????????, Michaelis--Menten ???????? ???(?????/?????)?? ????? ???????????.
???????$\cdot$???????????? ?????? ????????(identifiability) ???????? ???? ?????? ?????? ??????? $\kappa$\_d??
????????.

(3) ?????????? B\_k(basis-A, ????????????). B\_k = $\theta$\_fw + (1?$\theta$\_fw)$\cdot$(f\_prot$\cdot$K\_prot +
f\_PL$\cdot$K\_PL + f\_cw$\cdot$K\_cw). ?????????$\cdot$????????$\cdot$??????? ???????????? ?????? ?????????(Chen 2025
SSLM?? K\_PL, Zhou 2025 ???????? K\_prot ???)???? ??????, (1?$\theta$\_fw) ??????? ?????$\rightarrow$??????
???????? ?????????????.

(4) ???????? ?????. ?????? ????????? ???? ??????????? ???????? ?????????, ?? ?????? f\_xy(TSCF
???? ?????)?? root$\rightarrow$shoot ???????? ????????? ?????? lever??(??????????? ???? Casparian
???????? ?????????? ?????? ????? f\_xy?1). ????????? ????? ??????(L\_Ph)??? ?????, ????????? pH
???????????? ????????? ????????/???????? ??????? ?????. ???????????(????? ???????? $\phi$, ?????? ??????
?????? $\gamma$(t), ??? ???? ??????)??? ???????? ????????.

(5) ????? ??????. Freundlich/?????? Kd ?????????? ????? ????????? ???????? ??????? ????????????,
??????????? ????? HYDRUS-1D ??????(phydrus ?????)??? ??????? ??????????? C\_w\textasciicircum{}o(t)?? Q\_TP(t)???
???????????(????????? ???????? ????? ??? ??????, ????????? ???????? ??????).

(6) ??????????$\cdot$?????$\cdot$??????. ??????????? Tier 0--3???? ??????????(???????? ?????? / BAF??
?????? ???????? ????????? / ?????? ?????? ????? / QSPR$\cdot$??????). ????? Koc??? Higgins--Luthy
QSPR??, ???????? pKa?? ??????????, ????? ????????(f\_xy, L\_Ph)??? BAF ???????? ??????????.
????????? Yamazaki(2023, ????? ????? gradient), Tang(2026, ???????? TF), Kim(2019, ?????
BAF)?? ???? ??????????.

(7) ???? ??????$\cdot$???????. RDKit?? SMILES??? ???? ????????? ??????? ????????? read-across
????? QSPR?? ???????? ???????? ????????, Streamlit ?????????? ??????--????? ?????? colormap,
????????, ???????? ??????? ??????????.

Planned approaches:
\begin{itemize}
  \item four-compartment dynamic ODE (BDF)
  \item GHK + Michaelis-Menten hybrid root uptake
  \item basis-A binding factor from measured K\_PL / K\_prot
  \item f\_xy xylem loading + L\_Ph phloem grain feed
  \item calibration/validation vs Yamazaki 2023, Tang 2026, Kim 2019
\end{itemize}

\section{Hypotheses and findings}

\subsection{The four-compartment uptake ODE conserves PFAS mass: total compartment burden equals cumulative net root influx minus xylem/phloem exports and growth dilution, to solver tolerance.}
\begin{itemize}
  \item Hypothesis id: \texttt{hyp-mass} (confirmatory)
  \item Decision rule (threshold): max abs relative mass-balance residual over the season is < 1e-5 => support (mass-conserving to solver tolerance); >= 1e-5 => refute
  \item \textbf{Status: supported} --- confidence 9e-06 (credence)
  \item Basis: threshold rule: statistic 'point'=1e-06 < 1e-05 is met (combine='latest', margin=9e-06)
  \item \textit{Evidence validity:} referent=formal; data\_source(s)=generated --- in-silico/computational result
\end{itemize}

\subsection{An inside-negative root plasmalemma electrostatically excludes the PFAS anion (electrochemical factor e\textasciicircum{}N \textasciitilde{} 1e2), so passive diffusion alone cannot drive net uptake.}
\begin{itemize}
  \item Hypothesis id: \texttt{hyp-anion} (confirmatory)
  \item Decision rule (qualitative): GHK electrodiffusion with measured E\_m=-120 mV and z=-1 yields an anion-exclusion factor e\textasciicircum{}N of order 1e2 (the IOC-formulation signature) => support; no exclusion => refute
  \item \textbf{Status: supported} --- confidence strong (graded)
  \item Basis: qualitative judge: Both panelists agree the e\textasciicircum{}N\textasciitilde{}1e2 exclusion is a closed-form consequence of Tier-0 inputs (E\_m, z), independently asserted by test\_plant\_module; it is a structural/formal property, not a fitted one, so the IOC signature holds.
  \item \textit{Evidence validity:} referent=formal; data\_source(s)=generated --- in-silico/computational result
\end{itemize}

\subsection{Calibration to the Yamazaki (2023) dataset constitutes out-of-sample predictive validation of the model's tissue BAFs.}
\begin{itemize}
  \item Hypothesis id: \texttt{hyp-yamazaki} (confirmatory)
  \item Decision rule (qualitative): the Yamazaki agreement is OUT-OF-SAMPLE predictive validation (independent data not used to fit the per-congener transport parameters) => support; an in-sample / saturated reproduction => refute
  \item \textbf{Status: refuted} --- confidence strong (graded)
  \item Basis: qualitative judge: Decisive panel reasoning: the W2 transport parameters are FIT to the same Yamazaki tissues they then 'predict' (\textasciitilde{}3 params / 3 obs per congener), so RMSE 0.029 is saturated in-sample reproduction. The a-priori test (monotone f\_xy, --rec) gives the REAL predictive error log10 RMSE 0.837 (straw off 6-40x) -- the model does not predict out of sample. The predictive-validation hypothesis is refuted quantitatively.
  \item \textit{Evidence validity:} referent=empirical; data\_source(s)=measured
\end{itemize}

\subsection{The model predicts brown-rice grain PFAS accumulation accurately enough to support dietary risk assessment.}
\begin{itemize}
  \item Hypothesis id: \texttt{hyp-grain} (confirmatory)
  \item Decision rule (qualitative): model grain (brown-rice) BAF/TF matches measured within a factor adequate for dietary risk assessment => support; systematic structural under/over-prediction => refute
  \item \textbf{Status: refuted} --- confidence strong (graded)
  \item Basis: qualitative judge: Decisive panel reasoning: against MEASURED data the model under-predicts grain by \textasciitilde{}3-8x (Tang PFOA endosperm 0.11 vs 0.95) and only matches Kim by forcing a single-point L\_Ph anchor. A structural under-prediction of the risk-relevant compartment cannot support dietary risk assessment => refute.
  \item \textit{Evidence validity:} referent=empirical; data\_source(s)=measured
\end{itemize}

\subsection{Per-congener soil sorption (Koc chain-length QSPR) is strongly congener-dependent, so the one-way soil coupling cannot be replaced by a single constant pore-water concentration.}
\begin{itemize}
  \item Hypothesis id: \texttt{hyp-soil} (confirmatory)
  \item Decision rule (threshold): per-congener soil sorption (Koc chain-length QSPR) spans > 2 log10 units across the congener set, so a single constant pore-water Cwo cannot represent all congeners (the soil-coupling rationale) => support; < 2 log10 spread => refute
  \item \textbf{Status: supported} --- confidence 0.909 (credence)
  \item Basis: threshold rule: statistic 'point'=4.4 > 2 is met (combine='latest', margin=2.4)
  \item \textit{Evidence validity:} referent=formal; data\_source(s)=generated --- in-silico/computational result
\end{itemize}

\subsection{The SMILES structure front-end parameterizes a known PFAS by measured read-across, reproducing the curated congener record.}
\begin{itemize}
  \item Hypothesis id: \texttt{hyp-smiles} (confirmatory)
  \item Decision rule (qualitative): the SMILES front-end reproduces the curated measured-parameter model for a KNOWN PFAS via read-across (structure -> same Compound) => support; a mismatch => refute
  \item \textbf{Status: supported} --- confidence strong (graded)
  \item Basis: qualitative judge: Both panelists agree this is a faithful structure->parameter mapping, an independently tested software property (23/23). It is SUPPORTED only for KNOWN structures (read-across); novel-structure prediction is out of scope and remains provisional.
  \item \textit{Evidence validity:} referent=formal; data\_source(s)=generated --- in-silico/computational result
\end{itemize}

\section{Evidence}
\begin{itemize}
  \item \texttt{evi-literature} (literature): finding=\{"acquired": [\{"doi": "10.1016/j.jhazmat.2025.141017"\}, \{"doi": "10.1016/j.scitotenv.2019.03.240"\}, \{"doi": "10.1029/201
  \item \texttt{evi-mass} (experiment\_run): point=1e-06, finding=mass-balance residual bounded by rel 1e-6 / abs 1e-9; tests/test\_plant\_module.py::test\_mass\_conservation\_source\_is\_root\_
  \item \texttt{evi-anion} (experiment\_run): finding=N = z*E\_m*F/RT = +4.67 at E\_m=-120 mV, z=-1, T=298.15 K => e\textasciicircum{}N = 106.8 (test\_plant\_module asserts e\textasciicircum{}N \textasciitilde{} 107); the GHK te
  \item \texttt{evi-yamazaki} (experiment\_run): point=0.029, finding=reproduce\_demo.py: W2 transport fit reproduces Yamazaki root/straw/grain BAFs across 11 congeners at log10 RMSE 0.029 (e
  \item \texttt{evi-yamazaki-apriori} (experiment\_run): point=0.837, finding=reproduce\_demo.py --rec: with the theory/QSPR MONOTONE f\_xy (a-priori, NOT fit to the tissue BAFs) the predictive error 
  \item \texttt{evi-grain-tang} (experiment\_run): point=0.11, finding=Tang 2026 OOS: PFOA endosperm TF model 0.11 vs measured 0.95 (dw); grain structurally \textasciitilde{}3-8x under across congeners; not 
  \item \texttt{evi-grain-kim} (experiment\_run): finding=Kim 2019: grain BAF reproduced only by FORCING L\_Ph (0.07 -> 4.43, L\_Ph\textasciitilde{}0.84) as a single-point in-sample anchor; Kim is
  \item \texttt{evi-soil} (experiment\_run): point=4.4, finding=literature\_params.koc (Higgins-Luthy +0.55/CF2 QSPR): Koc(PFBA C4)=2.15 vs Koc(PFDoDA C12)=53985 L/kg => 4.40 log10 spre
  \item \texttt{evi-smiles} (experiment\_run): finding=tests/test\_pfas\_structure.py: 23 passed (RDKit). A canonical SMILES that matches a curated congener rebuilds the SAME Co
\end{itemize}

\section{Pending agent judgments}
These hypotheses use a non-numeric rule and await an in-session agent verdict (no autonomous LLM call):
\begin{itemize}
  \item \texttt{hyp-anion} (qualitative): GHK electrodiffusion with measured E\_m=-120 mV and z=-1 yields an anion-exclusion factor e\textasciicircum{}N of order 1e2 (the IOC-formulation signature) => support; no exclusion => refute
  \item finding: N = z*E\_m*F/RT = +4.67 at E\_m=-120 mV, z=-1, T=298.15 K => e\textasciicircum{}N = 106.8 (test\_plant\_module asserts e\textasciicircum{}N \textasciitilde{} 107); the GHK term in root\_uptake() excludes the anion accordingly.
  \item \texttt{hyp-yamazaki} (qualitative): the Yamazaki agreement is OUT-OF-SAMPLE predictive validation (independent data not used to fit the per-congener transport parameters) => support; an in-sample / saturated reproduction => refute
  \item finding: reproduce\_demo.py: W2 transport fit reproduces Yamazaki root/straw/grain BAFs across 11 congeners at log10 RMSE 0.029 (e.g. PFOA root 0.49/0.49, straw 0.83/0.83 pred/obs). The fit assigns \textasciitilde{}3 transport params per congener against 3 tissue observations => SATURATED: reproduction is structurally guaranteed and is NOT out-of-sample prediction (CLAUDE.md ?6).
reproduce\_demo.py --rec: with the theory/QSPR MONOTONE f\_xy (a-priori, NOT fit to the tissue BAFs) the predictive error is log10 RMSE 0.837 vs the saturated W2 fit's 0.029. Straw is off 6-40x (PFBA 45/11, PFBS 33/2.2). The model does NOT a-priori predict the Yamazaki tissue BAFs -- a quantitative confirmation of the REFUTED verdict.
  \item \texttt{hyp-grain} (qualitative): model grain (brown-rice) BAF/TF matches measured within a factor adequate for dietary risk assessment => support; systematic structural under/over-prediction => refute
  \item finding: Tang 2026 OOS: PFOA endosperm TF model 0.11 vs measured 0.95 (dw); grain structurally \textasciitilde{}3-8x under across congeners; not closable by L\_Ph / lipid tuning (docs/tang2026\_grain\_units\_exploration.md).
Kim 2019: grain BAF reproduced only by FORCING L\_Ph (0.07 -> 4.43, L\_Ph\textasciitilde{}0.84) as a single-point in-sample anchor; Kim is grain-only, so root/straw TF are unconstrained.
  \item \texttt{hyp-smiles} (qualitative): the SMILES front-end reproduces the curated measured-parameter model for a KNOWN PFAS via read-across (structure -> same Compound) => support; a mismatch => refute
  \item finding: tests/test\_pfas\_structure.py: 23 passed (RDKit). A canonical SMILES that matches a curated congener rebuilds the SAME Compound from params/parameters.json (measured read-across). CAVEAT: for NOVEL structures f\_xy is provisional (QSPR/interpolated), NOT validated --- this claim is scoped to known structures only.
\end{itemize}

\section{References}
\begin{itemize}
  \item \url{https://doi.org/10.1016/j.jhazmat.2025.141017}
  \item \url{https://doi.org/10.1016/j.scitotenv.2019.03.240}
  \item \url{https://doi.org/10.1029/2019WR025432}
  \item \url{https://doi.org/10.1021/acs.est.0c07420}
\end{itemize}

\end{document}